\chapter{Preface}

This is a book about doing mathematics with a computer.

It has two ambitions at once. The first is to be a \emph{user's manual}: a
practical, example-driven guide to Mathilda, a computer algebra system you can
download, read, and run for free. The second is to be a \emph{textbook on
computational mathematics}---an account of how genuine, sometimes advanced,
mathematics is actually carried out inside a computer algebra system: how a
machine differentiates and integrates, factors an integer or a polynomial,
solves an equation exactly, sums a series in closed form, or races through a
million-element array without ever leaving machine precision.

Those two goals reinforce each other. You cannot really understand a computer
algebra system without seeing the mathematics it embodies, and you cannot
appreciate the mathematics without watching a real system do it. So this book is
built around transcripts of Mathilda at work, and it treats every one of them as
a small experiment worth explaining.

\section*{Every example is real}

There is a promise behind every transcript in this book:

\begin{quote}
Nothing you see printed as \texttt{Out[\,]} was typed by a human. Each example is
run through the Mathilda binary at the moment the book is built, and the output
is inserted automatically.
\end{quote}

This is not a stylistic flourish; it is a design decision baked into how the book
is produced. The inputs live in small program files; a build step feeds them to
Mathilda and captures exactly what it prints; and the typesetting includes that
captured text verbatim. There is, quite literally, nowhere for a hand-written
answer to hide. If Mathilda's behaviour changes, the next build tells the truth
about it. The mechanics are described in \texttt{book/CONTEXT.md} for the curious
and for anyone who wants to contribute.

One consequence is worth stating plainly: Mathilda's output form is not always
the form a textbook would choose. Terms may be reordered, an exponential may be
written \texttt{E\^{}x} rather than $e^{x}$, and an indefinite integral may differ
from a table by a constant. Everything printed is nonetheless mathematically
correct, and part of learning a computer algebra system is learning to read what
it says.

\section*{Links to the reference}

Mathilda has hundreds of built-in functions, each with its own reference page in
the online documentation. Whenever this book mentions a built-in for the first
time in a discussion---\B{Integrate}, \B{Solve}, \B{FactorInteger}, and so
on---the name is a live hyperlink to that reference page. The book teaches; the
reference pages enumerate every option and corner case. Read this for the story
and follow the links for the details.

\section*{How to read this book}

You do not need Mathematica, a licence, or a network connection to follow along;
you need only a copy of Mathilda, which you can build from source. Type the
\texttt{In[\,]} lines yourself and you will see the same \texttt{Out[\,]}. Later
chapters assume the vocabulary of earlier ones, but the mathematical chapters in
Part~II are largely independent---dip into calculus or number theory as your
interest takes you.

Welcome. Let us compute.
